\section{Research Objectives and Contributions}
\label{sec:objectives}

\subsection{Research Objectives}
This dissertation pursues four research objectives:
\begin{enumerate}
	\item \textbf{Characterize commercial ReRAM write latency.} Identify and
	      experimentally validate the dominant structural factors that effect
	      write latency in commercial ReRAM crossbar memories, including crossbar
	      architecture, embedded ECCs, and device aging.
	\item \textbf{Recover hidden device architecture non-invasively.} Develop a
	      technique to recover the bit-exact ECC parity submatrix from commercial
	      ReRAM products using only write-access timing, without physical probing
	      or device modification.
	\item \textbf{Apply structural knowledge to entropy characterization.}
	      Demonstrate that ECC-aware analysis is necessary for rigorous
	      characterization of stochastic write latency, and construct entropy
	      sources whose min-entropy can be evaluated under NIST SP 800-90B
	      requirements.
	      % \item \textbf{Establish a foundation for exploitable data remanence.}
	      % Connect the timing channels characterized in this dissertation to the
	      % broader class of hardware security threats that exploit physical device
	      % state, and identify the conditions under which prior cell states in
	      % commercial ReRAM may leave recoverable analog traces.
\end{enumerate}

\subsection{Contributions}

This dissertation makes the following contributions:

\paragraph{Timing Side-Channel Characterization.}
We present the first systematic characterization of write latency in
commercial ReRAM crossbar memories, identifying an array-level timing
property called \emph{polarity serialization} and embedded ECC activity
as the dominant sources of structured timing variance. We introduce the
directional Hamming distance (DHD) as a metric for classifying write
transitions and demonstrate that ECC parity bit activity is large enough
to obscure the stochastic device physics that prior security and entropy
analyses assumed they were measuring. These findings are validated on
four samples each of two commercial RAMXeed ReRAM products.

\paragraph{BREW-RC: Bit-Exact ECC Recovery from Write Timing.}
We introduce BREW-RC, the first technique to recover the bit-exact ECC
parity submatrix of a commercial ReRAM product using only write-access
timing, without physical probing or device modification. BREW-RC
formulates ECC recovery as a SAT problem, deriving constraints from
observed write latency distributions and recovering the complete
generator matrix from two commercial products. The recovered matrix
enables accurate timing models in the presence of embedded ECCs and is
validated against results from destructive reverse engineering,
achieving a perfect match under all valid bit orderings.

\paragraph{ECC-Aware Entropy Harvesting.}
We demonstrate that structured variance in commercial ReRAM write
latency is primarily explained by ECC activity and the resulting DHD,
and that prior entropy claims for commercial ReRAM TRNGs are overstated
as a result. Applying ECC-aware analysis, we isolate the deterministic
ECC contribution from the underlying stochastic filament dynamics and
construct stationary, wear-leveled measurement sequences with
min-entropy rigorously estimated under NIST SP 800-90B. We further
characterize how device aging affects switching dynamics and entropy
yield over device lifetime. The resulting TRNG passes NIST SP 800-22 and
PractRand and achieves 17$\times$ throughput improvement over prior
ECC-unaware methods.

% \paragraph{Toward ReRAM Data Remanence.} We connect the timing channels
% characterized in this dissertation to FPGA
% Pentimenti~\cite{drewes2023pentimento}, a recently discovered class of
% hardware security threats that recover prior device data through Bias
% Temperature Instability effects on programmable routing. We identify
% the structural and stochastic components established by this work as
% necessary conditions for reasoning about whether prior cell states in
% commercial ReRAM leave recoverable analog traces, and discuss the open
% research challenges that remain.
